\documentclass[10pt,a4paper]{amsart}
\usepackage[margin=19mm]{geometry}
\usepackage{amsmath,amssymb,mathtools}
\usepackage[scr=rsfs,cal=euler]{mathalfa}
\usepackage{xfrac}
\usepackage[unicode,hidelinks,bookmarks=false]{hyperref}
\newcommand{\ie}{i.e.\ }
\newcommand{\cf}{cf.\ }
\newcommand{\resp}{respectively\ }
\newcommand{\Subsection}{\subsection}
\providecommand{\nobreakdash}{\nobreak\hbox{-}}
\setlength{\emergencystretch}{2em}
\multlinegap0pt
\usepackage{amssymb}
\usepackage{xcolor}
\usepackage{dsfont}
\usepackage[shortlabels]{enumitem}
\usepackage{mathtools}
\usepackage{float}
\newtheorem{lemma}{Lemma}
\newcommand{\module}[1]{\left|#1\right|}
\newcommand{\norm}[1]{\left\|#1\right\|}
\newcommand{\parentheses}[1]{\left(#1\right)}
\title{On average population levels for models with directed diffusion in heterogeneous environments}
\author{Andr\'e Rickes, Elena Braverman}
\date{Source: arXiv:2601.03473v2 (CC BY 4.0)}
\begin{document}
\maketitle

\noindent\emph{Source and licence.} On average population levels for models with directed diffusion in heterogeneous environments. Version arXiv:2601.03473v2 (CC BY 4.0), distributed by arXiv under the Creative Commons Attribution 4.0 International licence, http://creativecommons.org/licenses/by/4.0/, as recorded in the arXiv licence metadata for this exact version and shown on the abstract page https://arxiv.org/abs/2601.03473v2. Full licence notice: \texttt{licenses/arxiv-2601.03473-CC-BY-4.0.txt}.\par\smallskip

\noindent\textit{Editorial scope.} Counted unit: the article's lemma 1 (source0.tex lines 586--601). The article's complete original proof of it is delivered with it (source0.tex lines 603--667, one \texttt{proof} environment of 65 lines). No mathematics is written, repaired or re-derived: the statement and the proof are the article's own lines, in the article's own notation, and no retained line is substituted or re-flowed.  Retained uncounted as the source-local context the proof needs: lines 276--281; lines 346--351.  Nothing was omitted from any retained span, so the package carries no omission record.  No retained line contains a citation, so the wrapper carries no bibliography and no bibliography entry.  No retained line contains a figure environment, an image-inclusion command or drawing code, so no image is delivered and the package declares no asset dependency.  Editorial formatting only, in addition: an AMS \texttt{amsart} preamble that restates the article's own declarations and macros used by the retained lines, namely the environments \texttt{lemma} on separate counters, together with the standard packages those lines need.  The retained environments print in this document as Lemma 1 in order of appearance, whereas the article prints the same items with the numbering of its own full document.  The complete native source of the article as distributed in the arXiv e-print for this exact version is bundled unchanged as \texttt{sources/192227/source0.tex}.
\begin{equation}\label{intro2}
    \begin{cases}
  \displaystyle       d\Delta\parentheses{\frac{u}{P}}+ru\parentheses{1-\frac{u}{K}}=0, & x \in \Omega,  \\
  \displaystyle            \frac{\partial}{\partial n}\parentheses{\frac{u}{P}}=0,  & x \in \partial\Omega.
    \end{cases} 
\end{equation} 

\begin{lemma}\label{Prop_1_DeAngelis}
    As $d\to0^+$, we obtain
    \begin{equation*}
         u_d\to K\hspace{0.5cm} in \hspace{0.5cm} L^\infty(\overline{\Omega})\cap W^{1,2}(\Omega).
    \end{equation*}
\end{lemma}

\begin{lemma}%[Our idea]
 \label{Lemma 1}
 Suppose that (A) holds and let $u_d$ be a solution to \eqref{intro2}.
    \begin{enumerate}
        \item If $\displaystyle\int_\Omega\nabla\parentheses{\frac{K}{P}}\cdot\nabla \parentheses{\frac{1}{r}}\,dx<0$, then for $d>0$ small, 
        \begin{equation*}
            \int_\Omega u_d\,dx > \int_\Omega K\,dx.
        \end{equation*}

        \item If $\displaystyle\int_\Omega\nabla\parentheses{\frac{K}{P}}\cdot\nabla \parentheses{\frac{1}{r}}\,dx>0$ , then for $d>0$ small, 
        \begin{equation*}
            \int_\Omega u_d\,dx <\int_\Omega K\,dx.
        \end{equation*}

    \end{enumerate}
\end{lemma}

\begin{proof}
%    [Generalized version of Donald L. DeAngelis et al.'s paper]
        Let us prove item~1, as item 2 is proven analogously. Multiply the first equation in  \eqref{intro2} by $\frac{K}{ru_d}$,
        \begin{equation*}
            d\Delta\parentheses{\frac{u_d}{P}}\frac{K}{ru_d}+K-u_d=0,
        \end{equation*}
        integrate it and use that $\frac{\partial}{\partial n}(\frac{u_d}{P})=0$ on $\partial\Omega$ to obtain
        \begin{equation}\label{Eq.1 Feb14th}
            -d\int_\Omega\nabla\parentheses{\frac{u_d}{P}}\cdot\nabla\parentheses{\frac{K}{ru_d}}\,dx+\int_\Omega(K-u_d)\,dx=0.
        \end{equation}
        
        Next, let us prove that
        \begin{equation*}
            \int_\Omega\module{\nabla\parentheses{\frac{u_d}{P}}\cdot\nabla\parentheses{\frac{K}{ru_d}}-\nabla\parentheses{\frac{K}{P}}\cdot\nabla\parentheses{\frac{1}{r}}}\,dx\to0 \textrm{ as } d\to0^+. 
        \end{equation*}
        For that purpose, notice that 
        \begin{align*}
            &\int_\Omega\module{\nabla\parentheses{\frac{u_d}{P}}\cdot\nabla\parentheses{\frac{K}{ru_d}}-\nabla\parentheses{\frac{K}{P}}\cdot\nabla\parentheses{\frac{1}{r}}}\,dx\\
            \leq&\int_\Omega\module{\nabla\parentheses{\frac{u_d}{P}-\frac{K}{P}}\cdot\nabla\parentheses{\frac{K}{ru_d}} }\,dx+\int_\Omega\module{\nabla\parentheses{\frac{K}{P}}\cdot\nabla\parentheses{\frac{K}{ru_d} -\frac{1}{r}}  }\,dx\\
            \leq&\norm{\nabla\parentheses{\frac{u_d}{P}-\frac{K}{P}}}_{L^2(\Omega)}\norm{\nabla\parentheses{\frac{K}{ru_d}}}_{L^2(\Omega)}+\norm{\nabla\parentheses{\frac{K}{P}}}_{L^\infty(\Omega)}\norm{\nabla\parentheses{\frac{K}{ru_d}-\frac{1}{r}}}_{L^1(\Omega)}\\
            =&I+II.
        \end{align*}
    From Proposition \ref{Prop_1_DeAngelis} we get that $u_d\to K$ on $L^\infty(\Omega)\cap W^{1,2}(\Omega)$ as $d\to0$, therefore the first term of $I$ vanishes. Moreover,
    \begin{equation*}
        \nabla\parentheses{\frac{K}{ru_d}}=\nabla\parentheses{\frac{K}{r}}\frac{1}{u_d}-\frac{K}{ru_d^2}\nabla u_d,
    \end{equation*}
    which is bounded, as $u_d\to K$ implies that $\nabla u_d$ is bounded and $u_d$ is far from zero almost everywhere. Therefore, $I\to0$ when $d\to0$.


    To estimate $II$, notice that
	\begin{align}
		\nonumber&\norm{\nabla\parentheses{\frac{K}{r u_d}-\frac{1}{r}}}_{L^1(\Omega)}\\
		\nonumber=&\int_\Omega\module{\frac{\nabla K}{ru_d}-\frac{K\nabla r}{r^2u_d}-\frac{K}{ru_d^2}\nabla u_d+\frac{\nabla r}{r^2}}\,dx\\
		\nonumber\leq&\int_\Omega\module{\frac{\nabla K}{ru_d}-\frac{\nabla u_d}{ru_d}}\,dx+\int_\Omega\module{\frac{\nabla u_d}{ru_d}-\frac{K}{ru_d^2}\nabla u_d}\,dx+\int_\Omega \module{\frac{K\nabla r}{r^2 u_d}-\frac{\nabla r}{r^2}}\,dx\\
		=&\int_\Omega\module{\frac{1}{ru_d}}\module{\nabla K-\nabla u_d}\,dx+\int_\Omega\module{\frac{\nabla u_d}{ru_d}}\module{1-\frac{K}{u_d}}\,dx+\int_\Omega\module{\frac{\nabla r}{r^2}} \module{\frac{K}{u_d}-1}\,dx.\label{Eq.2 Prop principal}
	\end{align}	
According to Proposition \ref{Prop_1_DeAngelis}, as $d\to0^+$, $u_d\to K$ on both $L^\infty(\Omega)$ and $W^{1,2}(\Omega)$ norms, therefore,  $|\frac{1}{u_d}|$ and $|\nabla u_d|$ are bounded on $\Omega$,      and therefore, each integral on \eqref{Eq.2 Prop principal} vanishes as $d\to0^+$.
    
%    \begin{align*}
%        \norm{\nabla\parentheses{\frac{K}{r u_d}-\frac{1}{r}}}_{L^1(\Omega)}&=\int_\Omega\module{\frac{\nabla(\xi u_d)}{(\xi u_d)^2}-\frac{\nabla r}{r^2}}\,dx\\
%        &\leq\int_\Omega\frac{1}{(\xi u_d)^2}|\nabla(\xi u_d)-\nabla r|\,dx+\int_\Omega|\nabla r|\,\module{\frac{1}{(\xi u_d)^2}-\frac{1}{r^2}}\,dx\\
%        &=III+IV.
%    \end{align*}
%    Since
%    \begin{equation*}
%        \nabla(\xi u_d)-\nabla r=\nabla\parentheses{\frac{r}{K}(u_d-K)},
%    \end{equation*}
%    we have that $III\to0$ by Proposition \ref{Prop_1_DeAngelis}. And since
%    \begin{equation*}
%        \frac{1}{(\xi u_d)^2}-\frac{1}{r^2}=\frac{r^2-(\xi u_d)^2}{(\xi u_d)^2r^2}=\frac{1}{(\xi u_d)^2r^2}(r-\xi u_d)(r+\xi u_d)=\frac{1}{u_d^2r^2}(K-u_d)(K+u_d),
%    \end{equation*}
%    we have that $IV\to0$ by proposition \ref{Prop_1_DeAngelis}.

    Therefore, as $d\to0^+$,
    \begin{equation*}
        \int_\Omega\nabla\parentheses{\frac{u_d}{P}}\cdot\nabla\parentheses{\frac{K}{ru_d}}\,dx=\int_\Omega\nabla\parentheses{\frac{K}{P}}\cdot\nabla\parentheses{\frac{1}{r}}\,dx+h(d),
    \end{equation*}
    where $\lim_{d\to0^+}h(d)=0$. And from \eqref{Eq.1 Feb14th}, we get
    \begin{equation*}
        \int_\Omega(K-u_d)\,dx=d\parentheses{\int_\Omega\nabla\parentheses{\frac{K}{P}}\cdot\nabla\parentheses{\frac{1}{r}}\,dx+h(d)}
    \end{equation*}
    which is negative for small enough $d>0$, since $\int_\Omega\nabla\parentheses{\frac{K}{P}}\cdot\nabla\parentheses{\frac{1}{r}}\,dx<0$.
    

\end{proof}

\end{document}
